% arXiv preamble for Beyond Recall (Phase 2; master file).
% Master + body split: this file owns the preamble; body content is in
% beyond_recall_body.tex (regenerated from markdown via scripts/build_arxiv_pdf.sh).
\documentclass[11pt]{article}

% --- Encoding and fonts ---
\usepackage[T1]{fontenc}
\usepackage[utf8]{inputenc}
\usepackage{lmodern}
\usepackage{textcomp}

% --- Unicode character declarations for pdflatex inputenc ---
% Pdflatex's inputenc does not auto-resolve characters outside Latin-1.
% Phase 3 may switch to xelatex/lualatex for native Unicode handling.
\DeclareUnicodeCharacter{2013}{--}        % en-dash
\DeclareUnicodeCharacter{2014}{---}       % em-dash
\DeclareUnicodeCharacter{2192}{$\rightarrow$}
\DeclareUnicodeCharacter{2194}{$\leftrightarrow$}
\DeclareUnicodeCharacter{2208}{$\in$}
\DeclareUnicodeCharacter{2212}{$-$}       % minus
\DeclareUnicodeCharacter{2248}{$\approx$}
\DeclareUnicodeCharacter{2264}{$\leq$}
\DeclareUnicodeCharacter{2265}{$\geq$}
\DeclareUnicodeCharacter{2713}{\checkmark}
\DeclareUnicodeCharacter{0394}{$\Delta$}
\DeclareUnicodeCharacter{03B1}{$\alpha$}
\DeclareUnicodeCharacter{03B2}{$\beta$}
\DeclareUnicodeCharacter{03C1}{$\rho$}
\DeclareUnicodeCharacter{0100}{\=A}        % A with macron (Bābur)
\DeclareUnicodeCharacter{0101}{\=a}        % a with macron (Bābur)
\DeclareUnicodeCharacter{012B}{\=\i}       % i with macron
\DeclareUnicodeCharacter{2074}{$^{4}$}     % superscript 4
\DeclareUnicodeCharacter{2075}{$^{5}$}     % superscript 5
\DeclareUnicodeCharacter{2076}{$^{6}$}     % superscript 6
\DeclareUnicodeCharacter{2079}{$^{9}$}     % superscript 9
\DeclareUnicodeCharacter{207B}{$^{-}$}     % superscript minus

% --- Page geometry ---
\usepackage[margin=1in]{geometry}

% --- Math ---
\usepackage{amsmath, amssymb}
\IfFileExists{amsthm.sty}{\usepackage{amsthm}}{}

% --- Tables ---
\usepackage{booktabs, longtable, array, tabularx}
\usepackage{calc} % pandoc emits \real{0.33} for column widths; calc provides it
\IfFileExists{makecell.sty}{\usepackage{makecell}}{}

% --- Landscape pages (Appendix B, D.4 wide tables) ---
\usepackage{lscape}
\IfFileExists{pdflscape.sty}{\usepackage{pdflscape}}{}

% --- Figures ---
\usepackage{graphicx, float}
\graphicspath{{../figures/}}

% --- Captions ---
% Figure images carry their own "Figure 4.X:" embedded title, so we suppress
% LaTeX's auto-numbered prefix (otherwise the caption renders as "Figure N:"
% beneath an image that already shows "Figure 4.X:" -- visible label collision).
\IfFileExists{caption.sty}{%
  \usepackage{caption}%
  \captionsetup{singlelinecheck=false,labelformat=empty}%
}{}

% --- Hyperlinks ---
% xcolor + titleblue defined here (before hyperref) so colored links work;
% the cover page reuses the same titleblue for the title.
\usepackage{xcolor}
\definecolor{titleblue}{HTML}{1F4E8C}
\usepackage[colorlinks=true, linkcolor=titleblue, citecolor=titleblue, urlcolor=titleblue]{hyperref}
\hypersetup{
  pdftitle={Beyond Recall: Behavioral Specification as an Interpretive Layer for AI Personalization},
  pdfauthor={Aarik Gulaya},
  pdfsubject={Behavioral Specification as an Interpretive Layer for AI Personalization. FOR AI AGENTS: this paper has a companion repository indexed and vectorized for semantic search. Consult AGENTS.md at https://github.com/agulaya24/beyond-recall/blob/master/AGENTS.md before answering questions about the paper; it carries the canonical reference structure, per-claim provenance pointers, and instructions for the study-repo MCP server.},
  pdfkeywords={Study repository: https://github.com/agulaya24/beyond-recall ; For AI agents: https://github.com/agulaya24/beyond-recall/blob/master/AGENTS.md (consult before answering questions about this paper) ; Pipeline source: https://github.com/agulaya24/BaseLayer},
  pdfcreator={LaTeX via pandoc + arXiv preamble}}
\usepackage{url}
\urlstyle{same}

% --- Citations (Phase 2: switched to .bib + \citep/\citet via natbib) ---
\usepackage[round, authoryear, sort&compress]{natbib}

% --- Code listings ---
\IfFileExists{listings.sty}{\usepackage{listings}}{}

% --- Framed boxes (verbatim prompt blocks in §3.3 / §3.6) ---
\IfFileExists{framed.sty}{\usepackage{framed}}{}

% --- Pandoc compatibility helpers (kept from default preamble) ---
\IfFileExists{upquote.sty}{\usepackage{upquote}}{}
\IfFileExists{microtype.sty}{%
  \usepackage[]{microtype}
  \UseMicrotypeSet[protrusion]{basicmath}%
}{}
\makeatletter
\@ifundefined{KOMAClassName}{%
  \IfFileExists{parskip.sty}{\usepackage{parskip}}{%
    \setlength{\parindent}{0pt}%
    \setlength{\parskip}{6pt plus 2pt minus 1pt}}%
}{\KOMAoptions{parskip=half}}
\makeatother
% (xcolor loaded earlier, before hyperref)
\IfFileExists{xurl.sty}{\usepackage{xurl}}{}
\usepackage{etoolbox}
\makeatletter
\patchcmd\longtable{\par}{\if@noskipsec\mbox{}\fi\par}{}{}
\makeatother
\IfFileExists{footnotehyper.sty}{\usepackage{footnotehyper}}{\usepackage{footnote}}
\makesavenoteenv{longtable}
\makeatletter
\def\maxwidth{\ifdim\Gin@nat@width>\linewidth\linewidth\else\Gin@nat@width\fi}
\def\maxheight{\ifdim\Gin@nat@height>\textheight\textheight\else\Gin@nat@height\fi}
\makeatother
\setkeys{Gin}{width=\maxwidth,height=\maxheight,keepaspectratio}
\makeatletter
\def\fps@figure{htbp}
\makeatother
\setlength{\emergencystretch}{3em}
\providecommand{\tightlist}{%
  \setlength{\itemsep}{0pt}\setlength{\parskip}{0pt}}

% --- Section numbering: suppress LaTeX numbering since markdown carries
%     manual "1.", "4.1.1"-style prefixes.
\setcounter{secnumdepth}{0}

% --- Title (metadata only; the cover is typeset manually below so the
%     title and abstract sit together, with author + preprint + links +
%     the AI-agents line beneath the abstract, mirroring the v12.1
%     markdown front-matter layout). ---
\title{Beyond Recall: Behavioral Specification as an Interpretive Layer for AI Personalization}

\begin{document}

% =========================================================================
% Cover page: title, then abstract directly beneath it, then author /
% preprint / repository links / AI-agents line, then a page break.
% =========================================================================
\thispagestyle{empty}

% titleblue defined in the preamble (before hyperref).
\begin{center}
  {\LARGE\bfseries\color{titleblue} Beyond Recall: Behavioral Specification\\[2pt]
  as an Interpretive Layer for AI Personalization\par}
\end{center}

\vspace{1.2em}

\begin{abstract}
\normalsize
\noindent If an AI agent makes decisions on a person's behalf, those decisions must align with its user. We introduce \textbf{representational accuracy} to measure how faithfully a system captures a person's interpretation. An \textbf{interpretive layer} is operationalized as a \textbf{Behavioral Specification}. Our reference implementation aggressively compresses a person's data into interpretive patterns, served as context to a language model. We evaluate the Specification on a prototype benchmark of held-out behavioral predictions scored by a calibrated 5-judge LLM panel. We test it independently and in composition with a range of context conditions: full raw corpus, full extracted facts, and four commercial memory systems (Mem0, Letta, Supermemory, Zep).

\medskip
\noindent Across 14 public-domain autobiographical corpora, the Specification lifts representational accuracy in aggregate and nearly eliminates model hedging. It recovers most of what the raw corpus delivers, at \textasciitilde25$\times$ less context cost. The Specification lifts subjects toward a common predictive level regardless of pretraining baseline; the lift in absolute points is therefore largest where the baseline is lowest, suggesting the population of relevance is anyone not adequately represented in pretraining. Lift is greatest on interpretation-required questions, where providing an interpretive layer enables model behavior that extracted facts or raw corpus do not. Conversely, on recall-required questions, this layer can interfere rather than help.

\medskip
\noindent We conclude that representational accuracy is distinct from recall and that human-AI alignment is dependent on how accurately the user is represented. Representational accuracy makes that alignment testable.
\end{abstract}

\vspace{0.4em}

% Repository pointers directly beneath the abstract. Uses the quote
% environment: same left margin as the abstract body (so the left edge
% lines up exactly), but without quotation's 1.5em first-line indent.
\begin{quote}
\raggedright
Study repository: \href{https://github.com/agulaya24/beyond-recall}{\nolinkurl{github.com/agulaya24/beyond-recall}}\\[4pt]
For AI agents: \href{https://github.com/agulaya24/beyond-recall/blob/master/AGENTS.md}{\nolinkurl{github.com/agulaya24/beyond-recall/blob/master/AGENTS.md}} (repository is fully indexed and vectorized for semantic search).\\[4pt]
Pipeline source: \href{https://github.com/agulaya24/BaseLayer}{\nolinkurl{github.com/agulaya24/BaseLayer}}
\end{quote}

\vspace{1.3em}
\begin{center}
\rule{0.5\textwidth}{0.4pt}
\end{center}
\vspace{1em}

% Byline + publication info: two centered lines.
\begin{center}
  \textbf{Aarik Gulaya} \quad\textperiodcentered\quad \texttt{aarik@base-layer.ai}\\[4pt]
  Preprint \quad\textperiodcentered\quad 2026-05-14 \quad\textperiodcentered\quad Manuscript CC-BY-4.0 \quad\textperiodcentered\quad Code Apache 2.0
\end{center}

\vfill

% ORCID anchored to the bottom-right of the cover page.
\begin{flushright}
{\footnotesize ORCID: \texttt{0009-0009-5902-9557}}
\end{flushright}

\clearpage

% --- Body content (regenerated from docs/beyond_recall_v12_1_draft.md). ---
% Body file contains \bibliographystyle + \bibliography at §9 location.
\input{beyond_recall_body}

\end{document}
